%% paper/sections/discussion.tex
%% Draft only. Limitations get their own subsection, six specific items.
\label{sec:discussion}

Three pre-registered structural hypotheses were tested under a decision
rule and effect-size floor frozen before any test-split evaluation
(\S\ref{sec:pre-registration}), and all three came back null
(\S\ref{sec:results}). MELD's H1 and H2 are the underpowered-but-real flavor
of null --- positive margins, confidence intervals excluding zero, both
still short of the pre-registered bar; DAiSEE's H2 replication is the
indistinguishable-from-zero flavor, agreeing across item-paired and
seed-paired tests alike. Read together, the honest summary is not
``persistence doesn't matter'' or ``entity structure doesn't matter'' ---
this study's power does not license either claim --- but that the specific,
falsifiable version of both claims this paper pre-registered did not survive
contact with a locked test set, on either dataset. The identity-shuffle
control (\S\ref{sec:shuffle-control}) did the job it was designed to do: it
gave persistence a real chance to show a distinguishable effect against
entity-localized cropping and token budget, under conditions where the
architecture, features, and everything else were held fixed, and the effect
it isolated did not clear the bar set for it in advance.

The efficiency results (\S\ref{sec:efficiency}) are not a consolation
finding dressed up to offset the null hypotheses --- they are a separate,
narrower claim about where compute goes, and they hold regardless of how H1
and H2 came out. Entity-localized cropping cuts encoder cost by $4.53\times$
unconditionally; end-to-end cost relative to a grid baseline depends on
hardware placement in a way that spans a $0.23\times$--$7.49\times$ range,
not a single number; and the attention mechanism this paper's architecture
is built around turns out to be a rounding error in the compute budget
compared to the frozen encoder it sits on top of. If token-count-driven
efficiency work in this space wants to target the dominant cost term, that
term is per-crop encoding, not fusion attention.

\subsection{Limitations}
\label{sec:limitations}

\textbf{Frozen backbone, no fine-tuning.} Every condition in this paper,
including the grid baseline, consumes the same frozen DINOv2~\cite{oquab2023dinov2} features
(\S\ref{sec:method}). This is deliberate --- it is what makes the matched-
feature comparison possible --- but it means no result here speaks to
whether entity-persistent tokenization would change what an end-to-end
fine-tuned encoder learns, only to what a fixed, off-the-shelf visual
representation supports downstream.

\textbf{One architecture family.} EPT-Former is a standard multi-head-attention
transformer applied to a differently-defined token stream
(\S\ref{sec:method}); this paper does not test whether the same entity-
persistent tokens would behave differently under a state-space or
convolutional fusion mechanism, only under attention.

\textbf{MELD's identity recovery is imperfect.}
%% outputs/meld_phase1_2_summary/REPORT.md: 62.2% purity, 16.7% chance floor
Clustering purity against a 6-character reference set is 62.2\% (16.7\%
chance floor) --- well above chance but far from clean tracking. The
purity-stratified analysis (\S\ref{sec:results}) argues the pooled H1 null
is not simply an artifact of this, but it cannot rule out that better
identity recovery would move the result; it can only show the result is not
obviously monotone in purity as measured.

\textbf{DAiSEE's $E{=}1$ degeneracy.} DAiSEE's primary condition has exactly
one entity slot (\S\ref{sec:pre-registration}), so A2 and H1 do not exist for
it by construction. The DAiSEE replication tests H2 only; it says nothing
about persistence on single-subject video, because there is no second slot
for identity to persist across.

\textbf{GPU detection cost is unmeasured, not zero.} Every end-to-end cost
regime that involves detection uses CPU-measured detection latency
(\S\ref{sec:efficiency}); the GPU-only regime could not be measured due to a
CUDA-version mismatch in the available \texttt{onnxruntime-gpu} build, not
because GPU detection is expected to be free. A working GPU detection
measurement could shift the low end of the reported ratio range.

\textbf{OUC-CGE's effective sample size is unrecoverable.} The dataset-audit
connected-components analysis
%% outputs/phase3_5_audit/TRACK1_TRACK2_REPORT.md sec.4: "plausibly in the tens, not the thousands"
(\S\ref{sec:dataset-audit}) finds the corpus's effective number of
independent, exchangeable recording units is plausibly in the tens, not the
thousands implied by its 7700-clip count, but that no similarity threshold
recovers a trustworthy estimate of the true number. This is why OUC-CGE is
reported only as a case study, never as a results-table entry: not just its
split, but its usable sample size, is not something this paper's methods can
put a reliable number on.
